%=======================================================
% CHAPTER 1: LINEAR CODES
%=======================================================

\begin{frame}[plain]
  \begin{center}\bfseries\centering
    \includegraphics[scale=1]{assets/branding/opa0}
    \vspace{-5.75cm}\\
    {\adjustbox{scale=2.15}{\color{TealDarkCyan}\;II.\;Linear Codes}}
    \bigskip\bigskip
    {\ }\vspace{4cm}\\
    \bigskip
  \end{center}
\end{frame}


%%--- EXISTING FRAME RESTORED ---
%\begin{frame}{\textbf{II. Linear Codes}}{The ambient space}
%\small
%\defbox[Hamming space]{
%Fix a finite field $\F_q$. The set $\F_q^n$ is simultaneously a vector space and a metric space with
%\[
%  d(x,y)=\#\{i:x_i\ne y_i\}.
%\]
%}
%\begin{itemize}
%  \item The metric is translation invariant: $d(x+z,y+z)=d(x,y)$.
%  \item The spheres are combinatorial objects:
%  \[
%    |B_t(0)|=\sum_{i=0}^t \binom{n}{i}(q-1)^i.
%  \]
%  \item Coding theory asks for large subsets with large mutual distance.
%\end{itemize}
%\end{frame}

%\begin{frame}{\textbf{II. Linear Codes}}{Codes as metric subsets}
%\small
%\defbox[Code]{
%A code of length $n$ over $\F_q$ is a subset $\code\subseteq \F_q^n$. Its minimum distance is
%\[
%  d(\code)=\min\{d(x,y):x,y\in\code,\ x\ne y\}.
%\]
%}
%\begin{block}{Mathematician's viewpoint}
%The fundamental problem is an extremal problem in a finite metric space:
%\[
%  A_q(n,d)=\max\{|\code|:\code\subseteq\F_q^n,\ d(\code)\ge d\}.
%\]
%\end{block}
%\end{frame}

\begin{frame}{\textbf{II. Linear Codes}}{Linear Code}
\small
\defbox[Linear Encoding]{
\textbf{Definition.}\; A \textbf{linear encoding} over \(\F_q\) is a linear map
\[
\F_q^k \longrightarrow \F_q^n,
\qquad
\vect{x}\longmapsto \vect{c}.
\]}

\defbox[Linear code]{
\textbf{Definition.}\; A \textbf{linear $[n,k,d]_q$ code} is a $k$-dimensional subspace \[
\code\le \F_q^n
\] with minimum distance $d$. The elements of \(\code\) are called \textbf{codewords}.}
\end{frame}

%\begin{frame}{\textbf{II. Linear Codes}}{Weight}
%\small
%\defbox[Hamming weight]{
%For $x=(x_1,\ldots,x_n)\in\F_q^n$,
%\[
%  \wt(x)=|\supp(x)|,\qquad \supp(x)=\{i:x_i\ne0\}.
%\]
%}
%\begin{block}{Lemma}
%For a linear code $\code$,
%\[
%  d(\code)=\min_{0\ne c\in\code}\wt(c).
%\]
%\end{block}
%\textbf{Proof.} If $x,y\in\code$, then $x-y\in\code$ and $d(x,y)=\wt(x-y)$.
%\end{frame}

%--- Slide: Generator Matrices ---
\begin{frame}{\textbf{II. Linear Codes}}{Generator Matrices}
\small
\defbox[Generator matrix]{
\textbf{Definition.}\; A \textbf{generator matrix} for an $\nkcode{n}{k}{q}$ code is a matrix $\vect{G}\in\F_q^{k\times n}$ whose rows form a basis of $\code$. Encoding is the linear map \[
\F_q^k\longrightarrow \F_q^n,\qquad \vect{m}\longmapsto \vect{m}\vect{G}.
\]}
%\begin{itemize}
%  \item The code is the row space of $G$.
%  \item Different bases give different generator matrices for the same code.
%%  \item Replacing $G$ by $UG$, $U\in\operatorname{GL}_k(\F_q)$, changes coordinates on messages only.
%\end{itemize}
\vfill
\textbf{Example (The Hamming Code $[7,4,3]_2$).}\; 
%The $[7,4,3]_2$ code has $k=4$ message bits and $n=7$ codeword bits. 
Consider \[
\vect{G} = \begin{bmatrix}
	1&0&0&0&0&1&1\\
	0&1&0&0&1&0&1\\
	0&0&1&0&1&1&0\\
	0&0&0&1&1&1&1
\end{bmatrix}_{4\times 7}
\] Encode message $\vect{m} = (1,0,1,1)$:
\[
\vect{c} = \vect{m}\vect{G} = \begin{bmatrix}
	1 & 0 & 1 & 1
\end{bmatrix}
\cdot \vect{G} = \begin{bmatrix}
	1 & 0 & 1 & 1 & 0 & 1 & 0
\end{bmatrix}.
\]
%\begin{center}
%\begin{tikzpicture}[scale=0.82]
%	% Show the codeword as colored blocks
%	\foreach \i/\v/\type in {1/1/msg, 2/0/msg, 3/1/msg, 4/1/msg, 5/0/par, 6/1/par, 7/0/par} {
%		\pgfmathsetmacro{\xp}{(\i-1)*0.85}
%		\ifnum\i<5
%		\node[minimum size=0.78cm, draw=TealDarkCyan, thick, rounded corners=3pt,
%		fill=TealDarkCyan!15, font=\small\bfseries] at (\xp,0) {$\v$};
%		\else
%		\node[minimum size=0.78cm, draw=AccentCyan, thick, rounded corners=3pt,
%		fill=AccentCyan!18, font=\small\bfseries] at (\xp,0) {$\v$};
%		\fi
%		\node[font=\tiny, TealDarkCyan] at (\xp,-0.58) {$c_{\i}$};
%	}
%	\draw[TealDarkCyan,    decorate, decoration={brace,amplitude=5pt,raise=3pt}]
%	(-0.4,0.4) -- (2.95,0.4) node[midway,above=7pt,font=\scriptsize]{message ($k=4$)};
%	\draw[AccentCyan, decorate, decoration={brace,amplitude=5pt,raise=3pt}]
%	(3.4,0.4) -- (5.55,0.4) node[midway,above=7pt,font=\scriptsize]{parity ($r=3$)};
%\end{tikzpicture}
%\end{center}
%\begin{itemize}
%\item Rate $R = 4/7 \approx 57\%$ — sending 4 bits of info for every 7 bits transmitted
%\item The 3 parity bits are \emph{linear combinations} of the 4 message bits
%\end{itemize}
\end{frame}

%\begin{frame}{\textbf{II. Linear Codes}}{Equivalence of descriptions}
%\small
%\thmbox[Row-space equivalence]{
%Two full-rank matrices $G,G'\in\F_q^{k\times n}$ generate the same code if and only if
%\[
%  G'=UG
%\]
%for some $U\in\operatorname{GL}_k(\F_q)$.
%}
%\begin{itemize}
%  \item This is just change of basis in the subspace $\code$.
%  \item The code is the invariant object; the matrix is a coordinate choice.
%  \item Many computational algorithms are Gaussian elimination on this statement.
%\end{itemize}
%\end{frame}

\begin{frame}{\textbf{II. Linear Codes}}{Systematic Form}
\small
\defbox[Systematic Form]{
\textbf{Definition.}\;	A code is said to be in \textbf{systematic form} if it has a generator matrix of the form \[
	\mathbf{G}=[\mathbf{I}_k\mid \mathbf{A}]
\]}
\begin{itemize}
  \item The message appears visibly in the first $k$ coordinates.
  \item The remaining $n-k$ coordinates are parity symbols.
  \item Systematic form is convenient, but it is not an intrinsic property of the code.
\end{itemize}
\begin{center}
\begin{tikzpicture}[x=1cm,y=1cm]
	\node[process] (msg) at (-4.55,0) {$\vect{m}$\\[-1pt]\scriptsize message};
	\foreach \i/\b in {1/1,2/0,3/1,4/1}{
		\node[bit teal, right={-0.01cm of msg}, xshift={0.45*\i cm-0.2cm}] (m\i) at (-5.8+0.05*\i,-1) {$\b$};
	}
	
	\node[process blue, minimum width=2.35cm, minimum height=1.48cm] (G) at (0,0) {
		\large $\vect{G}=[\vect{I}_k\mid \vect{P}]$
%		\\[4pt] \scriptsize row space $=\mathcal C$
	};
	
	\node[process amber] (code) at (4.45,0) {$\vect{c}=\vect{m}\vect{G}$\\[-1pt]\scriptsize codeword};
	\foreach \i/\b/\sty in {
		1/1/teal,2/0/teal,3/1/teal,4/1/teal,
		5/0/amber,6/0/amber,7/1/amber
	}{
		\ifnum\i<5
		\node[bit teal] at (2.6+0.46*\i,-1) {$\b$};
		\else
		\node[bit amber] at (2.6+0.46*\i,-1) {$\b$};
		\fi
	}
	\draw[decorate, decoration={brace, amplitude=4pt}, TealDarkCyan!75, line width=0.7pt]
	(4.66,-1.35) -- (2.86,-1.35)
	node[midway, below=0.12cm, tiny label] {message};
	\draw[decorate, decoration={brace, amplitude=4pt}, AmberCode!85, line width=0.7pt]
	(6.06,-1.35) -- (4.66,-1.35)
	node[midway, below=0.12cm, tiny label] {parity};
	
	\draw[flow] (msg) -- (G)
	node[midway, above=0.12cm, tiny label] {linear map};
	\draw[flow blue] (G) -- (code)
	node[midway, above=0.12cm, tiny label] {redundancy};
\end{tikzpicture}
\end{center}
%\begin{block}{Caution}
%Coordinate permutations preserve distance, but they change the visible positions of message and parity symbols.
%\end{block}
\end{frame}

\begin{frame}[fragile]{\textbf{II. Linear Codes}}{Parity-check Matrices}
\small
\defbox[Parity-check Matrix]{
A \textbf{parity-check matrix} is a matrix $\vect{H}\in\F_q^{(n-k)\times n}$ such that \[
	\code=\ker(\vect{H})=\{\vect{c}\in\F_q^n:\vect{H}\vect{c}^T=0\}.
\]}
\begin{itemize}
  \item If $\vect{G}$ generates $\code$, then $\vect{G}\vect{H}^T=0$.
%  \item The rows of $\vect{H}$ generate the dual code $\code^\perp$.
%  \item The equation $\vect{H}\vect{c}^T=0$ is a system of homogeneous parity constraints.
\end{itemize}
\begin{center}
\begin{tikzpicture}
  \matrix[matrix of nodes,
nodes={bit blue},
column sep=-0.8pt,
row sep=-0.8pt,
%left delimiter={[},
%right delimiter={]},
] (H) at (-3.65,0.35) {
	0&1&1&1&1&0&0\\
	1&0&1&1&0&1&0\\
	1&1&0&1&0&0&1\\
%	1&0&1&0&1&0&1\\
%	0&1&1&0&0&1&1\\
%	0&0&0&1&1&1&1\\
};
\node[tiny label, above=0.18cm of H] {$\vect{H}$};

%\foreach \i/\b/\sty in {1/1/teal,2/0/teal,3/1/teal,4/1/teal,5/1/error,6/0/teal,7/1/teal}{
%	\ifnum\i=5
%	\node[bit error] (v\i) at (-1.1+0.46*\i,-1.45) {$\b$};
%	\else
%	\node[bit teal] (v\i) at (-1.1+0.46*\i,-1.45) {$\b$};
%	\fi
%}
%\node[tiny label, below=0.12cm of v4] {$\vect{v}$};
%\draw[ErrorRed, line width=0.9pt] (v5.north west) -- (v5.south east);
%\draw[ErrorRed, line width=0.9pt] (v5.south west) -- (v5.north east);
%\node[ErrorRed!85!black, font=\scriptsize\bfseries, below=0.12cm of v5] {error};

\node[process, minimum width=1.7cm] (mult) at (.5,0.35) {$\vect{H}\vect{c}^T$};
\matrix[matrix of nodes,
nodes={bit error},
row sep=-0.8pt,
left delimiter={(},
right delimiter={)}
] (s) at (3.4,0.35) {
	1\\0\\0\\
};
\node[tiny label, above=0.18cm of s] {$\vect{s}$};

%\node[process red, minimum width=2.55cm] (lookup) at (6,0.35)
%{lookup column\\[-1pt]\scriptsize $\vect{s}=\vect{H}_5$};

\draw[flow] (H.east) -- (mult.west);
\draw[flow] (mult.east) -- (s.west);
%\draw[flow red] (s.south) -- (lookup.north);
%\draw[flow red] (lookup.south) .. controls (6,-0.5) and (0.7,-0.5) .. (v5.north)
%%node[midway, below, tiny label] {flip coordinate $5$}
%;
\end{tikzpicture}
\end{center}
\end{frame}

\begin{frame}{\textbf{II. Linear Codes}}{Duality}
\small
\defbox[Dual code]{
\textbf{Definition.}\; The dual of $\code\le\F_q^n$ is
\[
  \code^\perp=\{\vect{y}\in\F_q^n:\langle \vect{x},\vect{y}\rangle=0\text{ for all }\vect{x}\in\code\}.
\]}
\textbf{Remark (Dimension identity).}\; $ \dim\code+\dim\code^\perp=n.$
%\begin{itemize}
%  \item A parity-check matrix for $\code$ is a generator matrix for $\code^\perp$.
%  \item Duality exchanges constraints and codewords.
%\end{itemize}
\end{frame}

%\begin{frame}{\textbf{II. Linear Codes}}{Distance from checks}
%\small
%\thmbox[Column criterion]{
%Let $H$ be a parity-check matrix of $\code$. Then $d(\code)$ is the least number of columns of $H$ that are linearly dependent.
%}
%\begin{proof}
%A vector $c$ of weight $s$ satisfies $Hc^T=0$ exactly when a linear combination of the $s$ corresponding columns is zero.
%\end{proof}
%\begin{itemize}
%  \item No zero column means $d\ge2$.
%  \item No proportional columns means $d\ge3$.
%  \item This criterion is the entrance to Hamming codes.
%\end{itemize}
%\end{frame}

\begin{frame}{\textbf{II. Linear Codes}}{Syndromes}
\small
\defbox[Syndrome]{
\textbf{Proposition.}\; For a received word $\vect{y}=\vect{c}+\vect{e}$, 
\[
\vect{H}\vect{y}^T=
\vect{H}\vect{c}^T+
\vect{H}\vect{e}^T
=
\vect{H}\vect{e}^T.
\]}
\begin{itemize}
%  \item The syndrome depends only on the error coset.
  \item $\vect{H}\vect{y}^T=0$ if and only if $\vec{y}\in\code$.
%  \item Decoding is the problem of choosing a likely error in each coset.
\end{itemize}
\end{frame}

%%--- NEW: Syndrome Decoding Walkthrough ---
%\begin{frame}{\textbf{II. Linear Codes}}{Syndrome Decoding: Step by Step}
%\small
%Using the $[7,4,3]$ Hamming code with
%$H = \begin{bmatrix}1&0&1&0&1&0&1\\0&1&1&0&0&1&1\\0&0&0&1&1&1&1\end{bmatrix}$.
%
%\medskip
%\begin{columns}[T,onlytextwidth]
%  \begin{column}{.52\textwidth}
%    \textbf{Step 1.} Transmitted $c = (1,0,1,1,0,1,0)$.
%
%    \textbf{Step 2.} Noise flips bit 5: $e = (0,0,0,0,1,0,0)$.
%
%    \textbf{Step 3.} Received $v = (1,0,1,1,\textcolor{AccentCyan}{\mathbf{1}},1,0)$.
%
%    \smallskip
%    \textbf{Step 4.} Compute syndrome $s = Hv^T$:
%    \[
%    s = \begin{pmatrix}1\\0\\1\end{pmatrix}.
%    \]
%
%    \textbf{Step 5.} Column 5 of $H$ is $(1,0,1)^T = s$.
%
%    \textbf{Step 6.} Flip bit 5. Corrected! $\hat{c} = c$. \checkmark
%  \end{column}
%  \begin{column}{.45\textwidth}
%    \begin{tikzpicture}[scale=0.80]
%      % Syndrome → column lookup table
%      \node[font=\scriptsize\bfseries, TealDarkCyan] at (1.5,3.4) {Syndrome lookup};
%      \foreach \i/\s in {1/{001},2/{010},3/{011},4/{100},5/{101},6/{110},7/{111}} {
%        \pgfmathsetmacro{\yy}{3.0-(\i-1)*0.38}
%        \ifnum\i=5
%          \node[draw=AccentCyan,fill=AccentCyan!18,rounded corners=2pt,
%                minimum width=2.5cm,minimum height=0.32cm,font=\scriptsize]
%            at (1.5,\yy) {\textbf{col \i} $\to s=(1,0,1)^T$ \checkmark};
%        \else
%          \node[draw=TealDarkCyan!30,rounded corners=2pt,
%                minimum width=2.5cm,minimum height=0.32cm,font=\scriptsize]
%            at (1.5,\yy) {col \i\ $\to$ column $\i$ of $H$};
%        \fi
%      }
%    \end{tikzpicture}
%
%    \smallskip
%    \begin{block}{}
%    \centering\footnotesize
%    The syndrome \emph{is} the error address.\\
%    Zero syndrome $\Rightarrow$ no error detected.
%    \end{block}
%  \end{column}
%\end{columns}
%\end{frame}
%
%\begin{frame}{\textbf{II. Linear Codes}}{Cosets}
%\small
%\defbox[Coset decomposition]{
%The ambient space is a disjoint union of cosets:
%\[
%  \F_q^n=\coprod_{a\in A}(a+\code),
%\]
%where $|A|=q^{n-k}$.
%}
%\begin{itemize}
%  \item All words in the same coset have the same syndrome.
%  \item Syndromes parameterize cosets when $H$ has full rank.
%  \item Decoding assigns to each coset a chosen representative of small weight.
%\end{itemize}
%\end{frame}
%
%%--- NEW: Coset Decomposition Visualized ---
%\begin{frame}{\textbf{II. Linear Codes}}{Coset Decomposition: Concrete Example}
%\small
%Using the $[3,1,3]$ repetition code $\code=\{000,111\}$ inside $\F_2^3$ (8 elements).
%
%\medskip
%\begin{center}
%\begin{tikzpicture}[
%  cw/.style={draw=TealDarkCyan, thick, rounded corners=3pt,
%             minimum width=1.05cm, minimum height=0.55cm, font=\scriptsize\bfseries,
%             fill=TealDarkCyan!15},
%  ldr/.style={draw=AccentCyan, thick, rounded corners=3pt,
%              minimum width=1.05cm, minimum height=0.55cm, font=\scriptsize\bfseries,
%              fill=AccentCyan!18}
%]
%  % Header row
%  \node[font=\scriptsize\bfseries, TealDarkCyan] at (-1.2,1.1) {Leader};
%  \node[font=\scriptsize\bfseries, TealDarkCyan] at ( 0.3,1.1) {Coset};
%  \node[font=\scriptsize\bfseries, TealDarkCyan] at ( 1.9,1.1) {Syndrome};
%  \node[font=\scriptsize\bfseries, TealDarkCyan] at ( 3.3,1.1) {Decoding};
%  \draw[TealDarkCyan!40, thick] (-1.75,0.85) -- (4.1,0.85);
%  % Coset 0: leader 000, elements 000+000=000 and 000+111=111
%  \node[cw] at (-1.2, 0.5) {$000$};
%  \node[cw] at ( 0.0, 0.5) {$\{000,$};
%  \node[cw] at ( 0.7, 0.5) {$111\}$};
%  \node[font=\scriptsize, TealDarkCyan] at (1.9,0.5) {$s=00$};
%  \node[font=\scriptsize, TealDarkCyan] at (3.3,0.5) {no error};
%  % Coset 1: leader 100
%  \node[ldr] at (-1.2, 0.0) {$100$};
%  \node[cw]  at ( 0.0, 0.0) {$\{100,$};
%  \node[cw]  at ( 0.7, 0.0) {$011\}$};
%  \node[font=\scriptsize, AccentCyan] at (1.9,0.0) {$s=11$};
%  \node[font=\scriptsize, AccentCyan] at (3.3,0.0) {flip bit 1};
%  % Coset 2: leader 010
%  \node[ldr] at (-1.2,-0.5) {$010$};
%  \node[cw]  at ( 0.0,-0.5) {$\{010,$};
%  \node[cw]  at ( 0.7,-0.5) {$101\}$};
%  \node[font=\scriptsize, AccentCyan] at (1.9,-0.5) {$s=10$};
%  \node[font=\scriptsize, AccentCyan] at (3.3,-0.5) {flip bit 2};
%  % Coset 3: leader 001
%  \node[ldr] at (-1.2,-1.0) {$001$};
%  \node[cw]  at ( 0.0,-1.0) {$\{001,$};
%  \node[cw]  at ( 0.7,-1.0) {$110\}$};
%  \node[font=\scriptsize, AccentCyan] at (1.9,-1.0) {$s=01$};
%  \node[font=\scriptsize, AccentCyan] at (3.3,-1.0) {flip bit 3};
%\end{tikzpicture}
%\end{center}
%
%\begin{columns}[T,onlytextwidth]
%  \begin{column}{.48\textwidth}
%    \textcolor{TealDarkCyan}{\textbf{Reading the table}}
%    \begin{itemize}
%      \item Each row is one coset of $\code$
%      \item The \textcolor{AccentCyan}{\textbf{leader}} is the most likely error pattern
%      \item The syndrome uniquely identifies the coset
%    \end{itemize}
%  \end{column}
%  \begin{column}{.48\textwidth}
%    \begin{block}{All 8 words are covered}
%    The 4 cosets are disjoint and together contain all $2^3=8$ elements of $\F_2^3$. Every received word is in exactly one coset.
%    \end{block}
%  \end{column}
%\end{columns}
%\end{frame}

%\begin{frame}{\textbf{II. Linear Codes}}{Standard arrays}
%\small
%\defbox[Coset leader]{
%A coset leader is an element of least Hamming weight in a coset $a+\code$.
%}
%\begin{itemize}
%  \item The standard array lists cosets by their leaders.
%  \item Nearest-neighbor decoding subtracts the coset leader.
%  \item The method is conceptually clean but exponential in $n-k$.
%\end{itemize}
%\begin{block}{Principle}
%Good algebraic families replace the standard array by structure.
%\end{block}
%\end{frame}
%
%%--- NEW: Nearest-Neighbor Decision Regions ---
%\begin{frame}{\textbf{II. Linear Codes}}{Nearest-Neighbor Decision Regions}
%\small
%Each coset's leader defines a \textbf{decision region}: when you receive a word, you correct to the nearest codeword.
%
%\medskip
%\begin{center}
%\begin{tikzpicture}[scale=0.85]
%  % Draw F_2^3 as 3-cube with colors by coset
%  % Coset 0 (code): 000 and 111 — TealDarkCyan
%  % Coset 1: 100 and 011 — AccentCyan!70
%  % Coset 2: 010 and 101 — AccentCyan!50
%  % Coset 3: 001 and 110 — AccentCyan!30
%  % Positions: 000=(0,0), 100=(2,0), 010=(1,0.9), 110=(3,0.9)
%  %            001=(0.5,1.6), 101=(2.5,1.6), 011=(1.5,2.5), 111=(3.5,2.5)
%  % Draw edges of 3-cube
%  \draw[TealDarkCyan!30,thick] (0,0)--(2,0); \draw[TealDarkCyan!30,thick] (0,0)--(1,0.9);
%  \draw[TealDarkCyan!30,thick] (0,0)--(0.5,1.6); \draw[TealDarkCyan!30,thick] (2,0)--(3,0.9);
%  \draw[TealDarkCyan!30,thick] (2,0)--(2.5,1.6); \draw[TealDarkCyan!30,thick] (1,0.9)--(3,0.9);
%  \draw[TealDarkCyan!30,thick] (1,0.9)--(1.5,2.5); \draw[TealDarkCyan!30,thick] (0.5,1.6)--(2.5,1.6);
%  \draw[TealDarkCyan!30,thick] (0.5,1.6)--(1.5,2.5); \draw[TealDarkCyan!30,thick] (3,0.9)--(3.5,2.5);
%  \draw[TealDarkCyan!30,thick] (2.5,1.6)--(3.5,2.5); \draw[TealDarkCyan!30,thick] (1.5,2.5)--(3.5,2.5);
%  % Nodes colored by coset
%  \node[circle,draw=TealDarkCyan,fill=TealDarkCyan!25,minimum size=0.62cm,font=\tiny\bfseries]
%    at (0,0) {$000$};
%  \node[circle,draw=AccentCyan,fill=AccentCyan!20,minimum size=0.62cm,font=\tiny\bfseries]
%    at (2,0) {$100$};
%  \node[circle,draw=AccentCyan,fill=AccentCyan!15,minimum size=0.62cm,font=\tiny\bfseries]
%    at (1,0.9) {$010$};
%  \node[circle,draw=TealDarkCyan!60,fill=TealDarkCyan!10,minimum size=0.62cm,font=\tiny\bfseries]
%    at (3,0.9) {$110$};
%  \node[circle,draw=TealDarkCyan!60,fill=TealDarkCyan!10,minimum size=0.62cm,font=\tiny\bfseries]
%    at (0.5,1.6) {$001$};
%  \node[circle,draw=AccentCyan,fill=AccentCyan!15,minimum size=0.62cm,font=\tiny\bfseries]
%    at (2.5,1.6) {$101$};
%  \node[circle,draw=AccentCyan,fill=AccentCyan!20,minimum size=0.62cm,font=\tiny\bfseries]
%    at (1.5,2.5) {$011$};
%  \node[circle,draw=TealDarkCyan,fill=TealDarkCyan!25,minimum size=0.62cm,font=\tiny\bfseries]
%    at (3.5,2.5) {$111$};
%  % Labels for codewords
%  \node[font=\tiny,TealDarkCyan,below=3pt] at (0,0)     {\textbf{codeword}};
%  \node[font=\tiny,TealDarkCyan,above=3pt] at (3.5,2.5) {\textbf{codeword}};
%  % Decision arrows
%  \draw[->,AccentCyan,thick,bend right=20] (2,0) to (0,0);
%  \draw[->,AccentCyan,thick,bend right=15] (1,0.9) to (0,0);
%  \draw[->,AccentCyan,thick,bend right=15] (0.5,1.6) to (0,0);
%  \draw[->,AccentCyan!80,thick,bend left=20] (1.5,2.5) to (3.5,2.5);
%  \draw[->,AccentCyan!80,thick,bend left=15] (2.5,1.6) to (3.5,2.5);
%  \draw[->,AccentCyan!80,thick,bend left=15] (3,0.9) to (3.5,2.5);
%\end{tikzpicture}
%\end{center}
%
%\vspace{-4pt}
%Each non-codeword points to its \textbf{nearest codeword} (distance 1 away). This is nearest-neighbor decoding applied to the $[3,1,3]$ code.
%\end{frame}

\begin{frame}{\textbf{II. Linear Codes}}{Correction and detection}
\small
\thmbox[Distance theorem]{
A code of minimum distance $d$ detects every error of weight at most $d-1$ and corrects every error of weight at most
\[
  t=\left\lfloor\frac{d-1}{2}\right\rfloor.
\]
}
\begin{itemize}
  \item Detection: no nonzero error of weight $<d$ is a codeword.
  \item Correction: balls of radius $t$ around distinct codewords are disjoint.
\end{itemize}
\end{frame}

%%--- NEW: The Packing Picture ---
%\begin{frame}{\textbf{II. Linear Codes}}{The Packing Picture}
%\small
%\begin{columns}[T,onlytextwidth]
%  \begin{column}{.50\textwidth}
%    Think of $\F_q^n$ as a big space with $q^n$ points. A code places $q^k$ \textbf{codewords} in this space, each surrounded by an \textbf{error-correction ball} of radius $t$.
%
%    \medskip
%    \begin{center}
%    \begin{tikzpicture}[scale=0.88]
%      % Background space
%      \fill[TealDarkCyan!6] (-2.8,-2.3) rectangle (2.8,2.3);
%      \node[font=\tiny, TealDarkCyan!60] at (2.2,2.0) {$\F_2^n$};
%      % Three codewords with correction balls
%      \fill[AccentCyan!25] (-1.4, 0.9) circle (0.75);
%      \fill[TealDarkCyan]  (-1.4, 0.9) circle (3pt)
%        node[above=3pt,font=\tiny]{$c_1$};
%      \fill[AccentCyan!25] ( 1.3, 0.9) circle (0.75);
%      \fill[TealDarkCyan]  ( 1.3, 0.9) circle (3pt)
%        node[above=3pt,font=\tiny]{$c_2$};
%      \fill[AccentCyan!25] ( 0.0,-1.1) circle (0.75);
%      \fill[TealDarkCyan]  ( 0.0,-1.1) circle (3pt)
%        node[above=3pt,font=\tiny]{$c_3$};
%      % Label balls
%      \node[font=\tiny, AccentCyan!70!TealDarkCyan] at (-2.3,-0.5)
%        {balls of};
%      \node[font=\tiny, AccentCyan!70!TealDarkCyan] at (-2.3,-0.8)
%        {radius $t$};
%      % Show disjoint
%      \draw[dashed, TealDarkCyan!50] (-2.8,0) -- (2.8,0);
%      \node[font=\tiny, TealDarkCyan!50] at (2.1,-0.15) {balls don't overlap};
%    \end{tikzpicture}
%    \end{center}
%  \end{column}
%  \begin{column}{.47\textwidth}
%    \thmbox[Hamming bound]{
%    $q^k \cdot |B_t| \le q^n$, where $|B_t| = \sum_{i=0}^{t}\binom{n}{i}(q-1)^i$.
%    }
%
%    \medskip
%    \textcolor{TealDarkCyan}{\textbf{Interpretation}}
%    \begin{itemize}
%      \item Each codeword "claims" all points within distance $t$
%      \item These claimed regions cannot overlap
%      \item More codewords $\Rightarrow$ smaller (or no) balls $\Rightarrow$ less correction
%    \end{itemize}
%
%    \medskip
%    \defbox[Perfect code]{
%    A code meeting the bound with equality tiles $\F_q^n$ exactly — every point is in exactly one ball. \textbf{No space is wasted.}
%    }
%  \end{column}
%\end{columns}
%\end{frame}

\begin{frame}{\textbf{II. Linear Codes}}{Packing}
\small
\thmbox[Hamming bound]{
If $\code$ is an $[n,k,d]_q$ code and $t=\lfloor(d-1)/2\rfloor$, then
\[
  q^k\sum_{i=0}^t \binom{n}{i}(q-1)^i\le q^n.
\]
}
%\begin{proof}
%The balls of radius $t$ around codewords are disjoint subsets of $\F_q^n$.
%\end{proof}
%\begin{block}{Perfect code}
%Equality means the radius-$t$ balls tile the whole Hamming space.
%\end{block}
\end{frame}

\begin{frame}{\textbf{II. Linear Codes}}{Singleton bound}
\small
\thmbox[Singleton bound]{
Every $[n,k,d]_q$ code satisfies
\[
  d\le n-k+1.
\]
}
%\begin{proof}
%Delete $d-1$ coordinates. Distinct codewords must remain distinct, so $q^k\le q^{n-d+1}$.
%\end{proof}
\begin{itemize}
  \item Codes meeting the bound are called MDS codes.
  \item Reed--Solomon codes are the standard nonbinary MDS examples.
\end{itemize}
\end{frame}

\begin{frame}{\textbf{II. Linear Codes}}{Gilbert--Varshamov idea}
\small
\thmbox[Existence principle]{
There exist linear codes whose parameters are good because a greedy construction cannot cover the whole space too early.
}
\[
  q^k\sum_{i=0}^{d-2}\binom{n-1}{i}(q-1)^i < q^n
  \quad\Rightarrow\quad
  \text{an }[n,k,d]_q\text{ code exists.}
\]
%\begin{itemize}
%  \item This is not a construction one wants to implement.
%  \item It proves that the asymptotic landscape is nontrivial.
%  \item Algebraic families try to approach such existence bounds explicitly.
%\end{itemize}
\end{frame}

\begin{frame}{\textbf{II. Linear Codes}}{Griesmer bound}
\small
\thmbox[Griesmer bound]{
For a linear $[n,k,d]_q$ code,
\[
  n\ge \sum_{i=0}^{k-1}\left\lceil \frac{d}{q^i}\right\rceil.
\]
}
%\begin{itemize}
%  \item This is stronger than Singleton in many small-parameter cases.
%  \item It comes from repeatedly puncturing on the support of a minimum word.
%  \item It is a linear-code bound: linearity is essential.
%\end{itemize}
\end{frame}

%\begin{frame}{\textbf{II. Linear Codes}}{Puncturing}
%\small
%\defbox[Puncturing]{
%Puncturing a code at coordinate $i$ means deleting the $i$-th coordinate from every word.
%}
%\begin{itemize}
%  \item Length decreases by one.
%  \item Dimension usually stays the same, but can drop.
%  \item Minimum distance can decrease by at most one.
%\end{itemize}
%\begin{block}{Use}
%Puncturing is the basic operation for inductive arguments on code parameters.
%\end{block}
%\end{frame}
%
%\begin{frame}{\textbf{II. Linear Codes}}{Shortening}
%\small
%\defbox[Shortening]{
%Shortening at coordinate $i$ means first selecting codewords with $c_i=0$, then puncturing at $i$.
%}
%\begin{itemize}
%  \item Length decreases by one.
%  \item Dimension can decrease by one.
%  \item Minimum distance does not decrease.
%\end{itemize}
%\begin{block}{Duality}
%Puncturing and shortening are dual operations:
%\[
%  (\text{puncture } \code)^\perp=\text{shorten }(\code^\perp).
%\]
%\end{block}
%\end{frame}

\begin{frame}{\textbf{II. Linear Codes}}{Weight enumerators}
\small
\defbox[Weight enumerator]{
The weight enumerator is
\[
  W_{\code}(X,Y)=\sum_{c\in\code}X^{n-\wt(c)}Y^{\wt(c)}.
\]
}
\begin{itemize}
  \item It records the distribution of weights.
  \item The minimum distance is the first nonzero exponent of $Y$ after the constant term.
  \item It is a coarser invariant than the code but a powerful one.
\end{itemize}
\end{frame}

%--- NEW: Weight Distribution Bar Chart [7,4,3] ---
\begin{frame}{\textbf{II. Linear Codes}}{Weight Distribution of the $[7,4,3]$ Code}
\small
The weight enumerator of $\operatorname{Ham}(3,2)$ is
$W(X,Y)=X^7+7X^4Y^3+7X^3Y^4+Y^7$.
This means the 16 codewords have weights:

\medskip
\begin{center}
\begin{tikzpicture}[scale=0.88]
  \draw[-Stealth,TealDarkCyan,thick] (-0.3,0) -- (8.5,0)
    node[right,font=\small]{weight};
  \draw[-Stealth,TealDarkCyan,thick] (0,-0.2) -- (0,3.2)
    node[above,font=\small]{count};
  % Gridlines
  \foreach \y in {1,2,3,4,5,6,7} {
    \draw[TealDarkCyan!15] (0,\y*0.35) -- (8.2,\y*0.35);
    \node[font=\tiny,TealDarkCyan!70,left=2pt] at (0,\y*0.35) {\y};
  }
  % Bars (weight 0..7, counts: 1,0,0,7,7,0,0,1)
  % weight 0 → count 1
  \fill[TealDarkCyan!40] (0.2,0) rectangle (0.8,0.35);
  \node[font=\tiny,TealDarkCyan] at (0.5,-0.18) {0};
  \node[font=\tiny,TealDarkCyan] at (0.5,0.43) {1};
  % weight 1 → count 0
  \node[font=\tiny,gray] at (1.5,-0.18) {1};
  % weight 2 → count 0
  \node[font=\tiny,gray] at (2.5,-0.18) {2};
  % weight 3 → count 7
  \fill[AccentCyan] (3.2,0) rectangle (3.8,2.45);
  \node[font=\tiny,TealDarkCyan] at (3.5,-0.18) {3};
  \node[font=\small\bfseries,AccentCyan] at (3.5,2.62) {7};
  % weight 4 → count 7
  \fill[AccentCyan] (4.2,0) rectangle (4.8,2.45);
  \node[font=\tiny,TealDarkCyan] at (4.5,-0.18) {4};
  \node[font=\small\bfseries,AccentCyan] at (4.5,2.62) {7};
  % weight 5 → 0
  \node[font=\tiny,gray] at (5.5,-0.18) {5};
  % weight 6 → 0
  \node[font=\tiny,gray] at (6.5,-0.18) {6};
  % weight 7 → 1
  \fill[TealDarkCyan!40] (7.2,0) rectangle (7.8,0.35);
  \node[font=\tiny,TealDarkCyan] at (7.5,-0.18) {7};
  \node[font=\tiny,TealDarkCyan] at (7.5,0.43) {1};
  % annotation
  \draw[<->,AccentCyan,thick] (3.2,2.7) -- (4.8,2.7)
    node[midway,above,font=\scriptsize]{$d=3$ (minimum)};
\end{tikzpicture}
\end{center}

\begin{columns}[T,onlytextwidth]
  \begin{column}{.48\textwidth}
    \textcolor{TealDarkCyan}{\textbf{What we learn}}
    \begin{itemize}
      \item No codeword has weight 1 or 2 $\Rightarrow d \ge 3$
      \item The 7 weight-3 words are the Fano plane lines
      \item Symmetric: $\code$ contains the all-ones word
    \end{itemize}
  \end{column}
  \begin{column}{.48\textwidth}
    \begin{block}{MacWilliams duality}
    The weight distribution of the dual (simplex) code follows by the MacWilliams transform applied to this bar chart.
    \end{block}
  \end{column}
\end{columns}
\end{frame}

\begin{frame}{\textbf{II. Linear Codes}}{MacWilliams identity}
\small
\thmbox[MacWilliams identity]{
For a linear $[n,k]_q$ code,
\[
  W_{\code^\perp}(X,Y)=\frac{1}{|\code|}W_{\code}(X+(q-1)Y,X-Y).
\]
}
\begin{itemize}
  \item The identity is a finite Fourier transform.
  \item It turns the weight distribution of $\code$ into that of $\code^\perp$.
  \item It is one of the central reasons duality is computationally useful.
\end{itemize}
\end{frame}

\begin{frame}{\textbf{II. Linear Codes}}{Self-orthogonality}
\small
\defbox[Self-orthogonal and self-dual]{
A code is self-orthogonal if $\code\subseteq\code^\perp$ and self-dual if $\code=\code^\perp$.
}
\begin{itemize}
  \item Self-dual binary codes have dimension $n/2$.
  \item Doubly-even self-dual codes have all weights divisible by $4$.
  \item Such codes connect coding theory to lattices and modular forms.
\end{itemize}
\end{frame}

\begin{frame}{\textbf{II. Linear Codes}}{Equivalent codes}
\small
\defbox[Monomial equivalence]{
Two $q$-ary linear codes are equivalent if one is obtained from the other by permuting coordinates and multiplying coordinates by nonzero scalars.
}
\begin{itemize}
  \item Over $\F_2$, this is just coordinate permutation.
  \item Equivalent codes have the same parameters and weight enumerator.
  \item Classification means classification up to this equivalence.
\end{itemize}
\end{frame}

\begin{frame}{\textbf{II. Linear Codes}}{Automorphisms}
\small
\defbox[Automorphism group]{
\[
  \Aut(\code)=\{\pi\in S_n:\pi(\code)=\code\}
\]
in the binary case.
}
\begin{itemize}
  \item A large automorphism group is a sign of hidden geometry.
  \item Symmetry reduces decoding complexity.
  \item Hamming and cyclic codes both arise from strong symmetry.
\end{itemize}
\end{frame}

\begin{frame}{\textbf{II. Linear Codes}}{Repetition code}
\small
\defbox[Repetition code]{
\[
  \code=\{(a,a,\ldots,a):a\in\F_q\}
\]
is an $[n,1,n]_q$ code.
}
\begin{itemize}
  \item It has maximal distance for dimension $1$.
  \item Its rate is $1/n$, so it is inefficient.
  \item It corrects $\lfloor(n-1)/2\rfloor$ errors by majority vote when $q=2$.
\end{itemize}
\end{frame}

\begin{frame}{\textbf{II. Linear Codes}}{Single parity-check code}
\small
\defbox[Parity-check code]{
\[
  \code=\{x\in\F_q^n:x_1+\cdots+x_n=0\}
\]
is an $[n,n-1,2]_q$ code.
}
\begin{itemize}
  \item The parity-check matrix is $H=[1\;1\;\cdots\;1]$.
  \item It detects one error but cannot correct one error.
  \item It is the dual of the repetition code.
\end{itemize}
\end{frame}

\begin{frame}{\textbf{II. Linear Codes}}{The simplex code}
\small
\defbox[Binary simplex code]{
Let $H_r$ be the matrix whose columns are all nonzero vectors of $\F_2^r$. The row space of $H_r$ is the simplex code.
}
\begin{itemize}
  \item It has parameters $[2^r-1,r,2^{r-1}]_2$.
  \item Every nonzero codeword has the same weight.
  \item Its dual is the binary Hamming code.
\end{itemize}
\end{frame}

\begin{frame}{\textbf{II. Linear Codes}}{Matroid viewpoint}
\small
\defbox[Vector matroid of a code]{
A parity-check matrix $H$ defines a matroid on the coordinate set $\{1,\ldots,n\}$ by linear dependence of columns.
}
\begin{itemize}
  \item Circuits of the matroid correspond to minimal support codewords.
  \item The minimum distance is the size of the smallest circuit.
  \item Projective geometry enters when the columns of $H$ are projective points.
\end{itemize}
\end{frame}

\begin{frame}{\textbf{II. Linear Codes}}{Tanner graphs}
\small
\defbox[Tanner graph]{
Given a parity-check matrix $H$, form a bipartite graph with variable nodes $1,\ldots,n$ and check nodes equal to the rows of $H$.
}
\begin{itemize}
  \item An edge joins variable $j$ to check $i$ when $H_{ij}\ne0$.
  \item Sparse $H$ gives iterative decoding algorithms.
  \item Classical algebraic codes often have dense checks; LDPC codes emphasize sparsity.
\end{itemize}
\end{frame}

%--- NEW: Tanner Graph for [7,4,3] ---
\begin{frame}{\textbf{II. Linear Codes}}{Tanner Graph for the $[7,4,3]$ Code}
\small
\begin{columns}[T,onlytextwidth]
  \begin{column}{.52\textwidth}
    The parity-check matrix $H_3$ gives a bipartite graph with $7$ \textcolor{TealDarkCyan}{\textbf{variable nodes}} and $3$ \textcolor{AccentCyan}{\textbf{check nodes}}.

    \medskip
    \begin{tikzpicture}[scale=0.78]
      % Variable nodes: v1..v7 at top
      \node[circle,draw=TealDarkCyan,fill=TealDarkCyan!15,minimum size=0.55cm,
            font=\tiny\bfseries] (v1) at (0,   2.0) {$v_1$};
      \node[circle,draw=TealDarkCyan,fill=TealDarkCyan!15,minimum size=0.55cm,
            font=\tiny\bfseries] (v2) at (0.9, 2.0) {$v_2$};
      \node[circle,draw=TealDarkCyan,fill=TealDarkCyan!15,minimum size=0.55cm,
            font=\tiny\bfseries] (v3) at (1.8, 2.0) {$v_3$};
      \node[circle,draw=TealDarkCyan,fill=TealDarkCyan!15,minimum size=0.55cm,
            font=\tiny\bfseries] (v4) at (2.7, 2.0) {$v_4$};
      \node[circle,draw=TealDarkCyan,fill=TealDarkCyan!15,minimum size=0.55cm,
            font=\tiny\bfseries] (v5) at (3.6, 2.0) {$v_5$};
      \node[circle,draw=TealDarkCyan,fill=TealDarkCyan!15,minimum size=0.55cm,
            font=\tiny\bfseries] (v6) at (4.5, 2.0) {$v_6$};
      \node[circle,draw=TealDarkCyan,fill=TealDarkCyan!15,minimum size=0.55cm,
            font=\tiny\bfseries] (v7) at (5.4, 2.0) {$v_7$};
      % Check nodes: c1..c3 at bottom
      \node[rectangle,draw=AccentCyan,fill=AccentCyan!18,minimum size=0.58cm,
            rounded corners=3pt,font=\tiny\bfseries] (c1) at (0.9, 0) {$c_1$};
      \node[rectangle,draw=AccentCyan,fill=AccentCyan!18,minimum size=0.58cm,
            rounded corners=3pt,font=\tiny\bfseries] (c2) at (2.7, 0) {$c_2$};
      \node[rectangle,draw=AccentCyan,fill=AccentCyan!18,minimum size=0.58cm,
            rounded corners=3pt,font=\tiny\bfseries] (c3) at (4.5, 0) {$c_3$};
      % Edges from H_3:
      % Row 1 (c1): positions 1,3,5,7
      \draw[TealDarkCyan!60] (v1)--(c1); \draw[TealDarkCyan!60] (v3)--(c1);
      \draw[TealDarkCyan!60] (v5)--(c1); \draw[TealDarkCyan!60] (v7)--(c1);
      % Row 2 (c2): positions 2,3,6,7
      \draw[AccentCyan!60] (v2)--(c2); \draw[AccentCyan!60] (v3)--(c2);
      \draw[AccentCyan!60] (v6)--(c2); \draw[AccentCyan!60] (v7)--(c2);
      % Row 3 (c3): positions 4,5,6,7
      \draw[TealDarkCyan!40] (v4)--(c3); \draw[TealDarkCyan!40] (v5)--(c3);
      \draw[TealDarkCyan!40] (v6)--(c3); \draw[TealDarkCyan!40] (v7)--(c3);
    \end{tikzpicture}
  \end{column}
  \begin{column}{.45\textwidth}
    \textcolor{TealDarkCyan}{\textbf{How to read the graph}}
    \begin{itemize}
      \item Variable node $v_j$ = bit $j$ of the codeword
      \item Check node $c_i$ = parity equation $i$
      \item Edge $(v_j, c_i)$ means bit $j$ appears in check $i$
    \end{itemize}

    \medskip
    \textcolor{TealDarkCyan}{\textbf{Each check node says:}}
    \[
      c_1{:}\; v_1 \oplus v_3 \oplus v_5 \oplus v_7 = 0
    \]
    \[
      c_2{:}\; v_2 \oplus v_3 \oplus v_6 \oplus v_7 = 0
    \]
    \[
      c_3{:}\; v_4 \oplus v_5 \oplus v_6 \oplus v_7 = 0
    \]

%    \medskip
%    \begin{block}{LDPC connection}
%    Modern LDPC codes use \emph{sparse} Tanner graphs --- most bits are in only 3--4 checks. This enables fast iterative decoding.
%    \end{block}
  \end{column}
\end{columns}
\end{frame}

\begin{frame}{\textbf{II. Linear Codes}}{Decoding hardness}
\small
\thmbox[Syndrome decoding problem]{
Given $H$, a syndrome $s$, and an integer $t$, decide whether there exists $e$ with
\[
  He^T=s,\qquad \wt(e)\le t.
\]
}
\begin{itemize}
  \item This problem is NP-complete in general.
  \item Code-based cryptography is built around this hardness.
  \item Algebraic code families are useful because they add trapdoor-like structure.
\end{itemize}
\end{frame}

\begin{frame}{\textbf{II. Linear Codes}}{Perfect codes}
\small
\defbox[Perfect code]{
A code correcting $t$ errors is perfect if the radius-$t$ balls around codewords partition $\F_q^n$.
}
\[
  |\code|\sum_{i=0}^t \binom{n}{i}(q-1)^i=q^n.
\]
\begin{itemize}
  \item Perfect codes are extremely rare.
  \item The binary Hamming codes are perfect one-error-correcting codes.
  \item The Golay codes are the exceptional perfect codes.
\end{itemize}
\end{frame}

%%--- NEW: Linear Codes in Practice ---
%\begin{frame}{\textbf{II. Linear Codes}}{Linear Codes in Practice}
%\small
%\begin{columns}[T,onlytextwidth]
%  \begin{column}{.50\textwidth}
%    \textcolor{TealDarkCyan}{\textbf{Where you use them every day}}
%
%    \medskip
%    \begin{tabular}{p{2.0cm}|p{3.2cm}}
%      \textbf{Application} & \textbf{Code} \\\hline
%      \rowcolor{TealDarkCyan!8}
%      ECC RAM & Extended Hamming $[72,64,4]_2$ \\
%      RAID-6 disks & Reed--Solomon over $\F_{2^8}$ \\
%      \rowcolor{TealDarkCyan!8}
%      USB / Ethernet & CRC (cyclic) \\
%      QR codes & Reed--Solomon RS$(255,\cdot)$ \\
%      \rowcolor{TealDarkCyan!8}
%      Satellite TV & Turbo / LDPC codes \\
%    \end{tabular}
%
%    \medskip
%    \textcolor{TealDarkCyan}{\textbf{Why linear specifically?}}
%    \begin{itemize}
%      \item Encoding $m \mapsto mG$ costs $O(nk)$ — fast
%      \item Syndrome $Hv^T$ costs $O(n(n{-}k))$ — fast
%      \item Parameters can be \emph{proved}, not just measured
%    \end{itemize}
%  \end{column}
%  \begin{column}{.47\textwidth}
%    \begin{tikzpicture}[scale=0.88]
%      % Bar chart: rate vs correction
%      \draw[-Stealth,TealDarkCyan,thick] (0,0)--(4.0,0)
%        node[right,font=\scriptsize]{Rate $R$};
%      \draw[-Stealth,TealDarkCyan,thick] (0,0)--(0,3.2)
%        node[above,font=\scriptsize]{$t$ errors corrected};
%      % Repetition [3,1]
%      \fill[TealDarkCyan!40] (0,0) rectangle (0.4,1.0);
%      \node[font=\tiny,rotate=90,white] at (0.2,0.5){Rep[3]};
%      % Hamming [7,4]
%      \fill[TealDarkCyan!60] (0.6,0) rectangle (1.2,0.5);
%      \node[font=\tiny,rotate=90,white] at (0.9,0.25){Ham[7,4]};
%      % BCH [15,7]
%      \fill[AccentCyan!70] (1.4,0) rectangle (2.0,1.0);
%      \node[font=\tiny,rotate=90,white] at (1.7,0.5){BCH[15,7]};
%      % RS [255,223]
%      \fill[AccentCyan] (2.2,0) rectangle (3.6,1.5);
%      \node[font=\tiny,rotate=90,white] at (2.9,0.75){RS[255,223]};
%      % Axis labels
%      \node[font=\tiny,TealDarkCyan] at (0.2,-0.25) {$\frac{1}{3}$};
%      \node[font=\tiny,TealDarkCyan] at (0.9,-0.25) {$\frac{4}{7}$};
%      \node[font=\tiny,TealDarkCyan] at (1.7,-0.25) {$\frac{7}{15}$};
%      \node[font=\tiny,TealDarkCyan] at (2.9,-0.25) {$\frac{223}{255}$};
%      \node[font=\tiny,TealDarkCyan] at (-0.3,0.5) {$1$};
%      \node[font=\tiny,TealDarkCyan] at (-0.3,1.0) {$2$};
%      \node[font=\tiny,TealDarkCyan] at (-0.3,1.5) {$16$};
%    \end{tikzpicture}
%  \end{column}
%\end{columns}
%\end{frame}

\begin{frame}{\textbf{II. Linear Codes}}{Why linear codes are central}
\small
\begin{columns}[T,onlytextwidth]
  \begin{column}{.48\textwidth}
    \defbox[Algebra]{
    \begin{itemize}\small
      \item subspaces
      \item duality
      \item matrices
      \item finite fields
    \end{itemize}
    }
  \end{column}
  \begin{column}{.48\textwidth}
    \defbox[Algorithms]{
    \begin{itemize}\small
      \item encoding by $G$
      \item checking by $H$
      \item syndrome decoding
      \item structured families
    \end{itemize}
    }
  \end{column}
\end{columns}
\end{frame}

\begin{frame}{\textbf{II. Linear Codes}}{The road to Hamming}
\small
\begin{block}{Question}
Can we construct a linear code that corrects one error and wastes no syndrome?
\end{block}
\begin{itemize}
  \item One error at one of $n$ positions gives $n$ possibilities.
  \item No error gives one more possibility.
  \item A parity-check matrix with $r$ rows has $2^r$ binary syndromes.
\end{itemize}
\[
  n+1=2^r
\]
is the numerology of binary Hamming codes.
\end{frame}
